\documentclass[a4paper]{article}
\usepackage[utf8]{inputenc}
\usepackage[T2A]{fontenc}
\usepackage[english,russian]{babel}
\usepackage[left=25mm, top=20mm, right=25mm, bottom=30mm, nohead, nofoot]{geometry}
\usepackage{amsmath,amsfonts,amssymb}
\usepackage{fancybox,fancyhdr}
\usepackage{enumitem}
\usepackage{listings}
\usepackage{xcolor}
\usepackage{hyperref}

\hypersetup{colorlinks=true, allcolors=[RGB]{010 090 200}}

\pagestyle{fancy}
\fancyhf{}
\fancyhead[R]{Информатика ЕГЭ — Задание 17}
\fancyfoot[R]{\thepage}
\headsep=10mm

\definecolor{easycolor}{RGB}{0,150,0}
\definecolor{medcolor}{RGB}{200,100,0}
\definecolor{hardcolor}{RGB}{180,0,0}

\newcommand{\easytag}{\textcolor{easycolor}{\textbf{[easy]}}}
\newcommand{\medtag}{\textcolor{medcolor}{\textbf{[medium]}}}
\newcommand{\hardtag}{\textcolor{hardcolor}{\textbf{[hard]}}}
\newcommand{\answer}[1]{\par\smallskip\textit{\textbf{Ответ:} #1}}
\newcommand{\solution}[1]{\par\textit{\textbf{Решение:} #1}}
\newcommand{\filehint}[1]{\par\smallskip\textbf{Файл с данными:} \texttt{#1} (скачать с \href{https://inf-ege.sdamgia.ru/}{СдамГИА})}

% ============================================================
% КАК ДОБАВИТЬ НОВУЮ ЗАДАЧУ:
% \item % ID: cs_ege_17_???_00N | \???tag | src: №XXXXX
% Условие...
% \filehint{имя_файла.txt}
% \answer{...}
% \solution{...}
% ============================================================

\begin{document}

\Large
\section*{Задание 17 — Работа с файлом данных}

\normalsize
\textit{Тема:} написание программы на Python для обработки числовых последовательностей из файла.\\
\textit{Источник:} \href{https://inf-ege.sdamgia.ru/}{СдамГИА — Информатика ЕГЭ}

\bigskip

% ========================
%  ЛЁГКИЕ ЗАДАЧИ (easy)
% ========================

\subsection*{\textcolor{easycolor}{Лёгкий уровень}}

\begin{enumerate}[label=\textbf{Задача \arabic*.}, leftmargin=*, itemsep=10mm]

  \item % ID: cs_ege_17_easy_001 | \easytag | src: №37359
  \easytag\ В файле содержится последовательность из 10\,000 целых положительных чисел.
  Каждое число не превышает 10\,000.
  Определите и запишите в ответе сначала \textbf{количество пар} элементов
  последовательности, у которых \textbf{сумма элементов кратна 117},
  затем \textbf{максимальную} из сумм элементов таких пар.
  В данной задаче под парой подразумевается два \emph{различных} элемента
  последовательности. Порядок элементов в паре не важен.
  \filehint{17.txt}
  \answer{Два числа: количество пар и максимальная сумма}
  \solution{
  Читаем файл, перебираем все пары $i < j$, проверяем кратность суммы числу 117.
  Оптимизация: группируем числа по остатку от деления на 117.}

  \bigskip

  \item % ID: cs_ege_17_easy_002 | \easytag | src: №37340
  \easytag\ В файле содержится последовательность из 10\,000 целых положительных чисел.
  Каждое число не превышает 10\,000.
  Определите и запишите в ответе сначала \textbf{количество пар} элементов
  последовательности, \textbf{разность которых чётна и хотя бы одно из чисел
  делится на 31}, затем \textbf{максимальную} из сумм элементов таких пар.
  В данной задаче под парой подразумевается два \emph{различных} элемента
  последовательности. Порядок элементов в паре не важен.
  \filehint{17.txt}
  \answer{Два числа: количество пар и максимальная сумма}
  \solution{Разность чётна $\Leftrightarrow$ числа одной чётности. Фильтруем пары по этому условию и по делимости хотя бы одного на 31.}

  \bigskip

  \item % ID: cs_ege_17_easy_003 | \easytag | src: №37370
  \easytag\ В файле содержится последовательность из 10\,000 целых положительных чисел.
  Каждое число не превышает 10\,000.
  Определите и запишите в ответе сначала \textbf{количество пар} элементов
  последовательности, у которых \textbf{разность элементов кратна 60 и хотя бы
  один из элементов кратен 15}, затем \textbf{максимальную} из разностей элементов
  таких пар. Под парой подразумевается два \emph{различных} элемента.
  Порядок не важен.
  \filehint{17.txt}
  \answer{Два числа: количество пар и максимальная разность}
  \solution{Разность кратна 60: числа дают одинаковый остаток от деления на 60. Дополнительно хотя бы одно кратно 15. Перебираем пары с соответствующей проверкой.}

  \bigskip

  \item % ID: cs_ege_17_easy_004 | \easytag | src: №37350
  \easytag\ В файле содержится последовательность из 10\,000 целых положительных чисел.
  Каждое число не превышает 10\,000.
  Определите и запишите в ответе сначала \textbf{количество пар} элементов
  последовательности, у которых \textbf{сумма нечётна, а произведение делится на 3},
  затем \textbf{максимальную} из сумм элементов таких пар.
  Под парой подразумевается два \emph{различных} элемента. Порядок не важен.
  \filehint{17.txt}
  \answer{Два числа: количество пар и максимальная сумма}
  \solution{Сумма нечётна $\Leftrightarrow$ одно число чётное, другое нечётное. Произведение делится на 3 $\Leftrightarrow$ хотя бы одно кратно 3.}

  \bigskip

  \item % ID: cs_ege_17_easy_005 | \easytag | src: №40992
  \easytag\ Файл содержит последовательность неотрицательных целых чисел, не превышающих
  10\,000. Назовём парой два \emph{идущих подряд} элемента последовательности.
  Определите количество пар, в которых \textbf{хотя бы один из двух элементов делится
  на 5} и \textbf{хотя бы один из двух элементов меньше среднего арифметического всех
  нечётных элементов} последовательности. В ответе запишите два числа: сначала
  количество найденных пар, а затем — максимальную сумму элементов таких пар.

  \textit{Например:} в последовательности $(8\ 10\ 2\ 9\ 5)$ есть две подходящие
  пары: $(10, 2)$ и $(9, 5)$, ответ: $2$ и $14$.
  \filehint{17.txt}
  \answer{Два числа: количество пар и максимальная сумма}
  \solution{Предварительно вычисляем среднее нечётных. Затем перебираем соседние пары.}

\end{enumerate}

\bigskip

% ========================
%  СРЕДНИЕ ЗАДАЧИ (med)
% ========================

\subsection*{\textcolor{medcolor}{Средний уровень}}

\begin{enumerate}[label=\textbf{Задача \arabic*.}, leftmargin=*, itemsep=10mm]

  \item % ID: cs_ege_17_med_001 | \medtag | src: №47014
  \medtag\ Файл содержит последовательность неотрицательных целых чисел, не превышающих
  10\,000. Назовём парой два \emph{идущих подряд} элемента последовательности.
  Определите количество пар, в которых \textbf{один из двух элементов делится на 5},
  а \textbf{другой меньше среднего арифметического всех нечётных элементов}
  последовательности. В ответе запишите два числа: сначала количество найденных
  пар, а затем — максимальную сумму элементов таких пар.

  \textit{Например:} в последовательности $(8\ 10\ 2\ 7\ 5\ 1)$ есть две подходящие
  пары: $(10, 2)$ и $(5, 1)$, ответ: $2$ и $12$.
  \filehint{17.txt}
  \answer{Два числа}
  \solution{Отличие от easy\_005: не просто «хотя бы один», а ровно один делится на 5, а другой меньше среднего.}

  \bigskip

  \item % ID: cs_ege_17_med_002 | \medtag | src: №47221
  \medtag\ В файле содержится последовательность целых чисел.
  Элементы могут принимать значения от $-10\,000$ до $10\,000$.
  Определите количество пар последовательности, в которых \textbf{только одно
  число оканчивается на 3}, а \textbf{сумма квадратов элементов пары не меньше
  квадрата максимального элемента последовательности, оканчивающегося на 3}.
  В ответе запишите два числа: сначала количество найденных пар, затем
  максимальную из сумм квадратов элементов таких пар.
  Под парой подразумевается два \emph{идущих подряд} элемента.
  \filehint{17.txt}
  \answer{Два числа}
  \solution{Предварительно находим максимальный элемент, оканчивающийся на 3 ($M$). Перебираем соседние пары: ровно одно из двух оканчивается на 3, $a^2+b^2 \geq M^2$.}

  \bigskip

  \item % ID: cs_ege_17_med_003 | \medtag | src: №48465
  \medtag\ Файл содержит последовательность целых чисел, по модулю не превышающих
  10\,000. Назовём парой два \emph{идущих подряд} элемента последовательности.
  Определите количество таких пар, в которых \textbf{запись ровно одного элемента
  заканчивается цифрой 6}, а \textbf{сумма квадратов элементов пары меньше, чем
  квадрат наименьшего из элементов последовательности, запись которых заканчивается
  цифрой 6}.
  В ответе запишите два числа: сначала количество найденных пар, затем
  максимальную сумму квадратов элементов этих пар.
  \filehint{17.txt}
  \answer{Два числа}
  \solution{Находим минимальный элемент, оканчивающийся на 6 ($m$). Условие: ровно один оканчивается на 6, $a^2+b^2 < m^2$.}

  \bigskip

  \item % ID: cs_ege_17_med_004 | \medtag | src: №55813
  \medtag\ В файле содержится последовательность целых чисел.
  Элементы могут принимать целые значения от $1$ до $100\,000$ включительно.
  Определите количество пар последовательности, в которых \textbf{только одно
  число трёхзначное}, и \textbf{сумма элементов пары кратна минимальному
  трёхзначному значению последовательности, оканчивающемуся на 5}.
  В ответе запишите два числа: сначала количество найденных пар, затем
  минимальную из сумм элементов таких пар.
  Под парой подразумевается два \emph{идущих подряд} элемента.
  \filehint{17.txt}
  \answer{Два числа}
  \solution{Предварительно находим минимальное трёхзначное число, оканчивающееся на 5 ($d$). Условие: ровно одно из двух — трёхзначное, $(a+b)\ \%\ d = 0$.}

  \bigskip

  \item % ID: cs_ege_17_med_005 | \medtag | src: №75254
  \medtag\ Файл содержит последовательность натуральных чисел, не превышающих $100\,000$.
  Назовём тройкой три \emph{идущих подряд} элемента последовательности.
  Определите количество троек, для которых выполняются следующие условия:
  \begin{itemize}
    \item в тройке есть четырёхзначные числа;
    \item в тройке \textbf{не более одного} числа, у которого остаток от деления
    на 5 равен остатку от деления на 5 \emph{минимального} элемента всей
    последовательности;
    \item в тройке \textbf{не менее двух} чисел, у которых остаток от деления на 7
    равен остатку от деления на 7 \emph{максимального} элемента всей
    последовательности.
  \end{itemize}
  В ответе запишите два числа: сначала количество найденных троек, затем
  максимальную величину суммы элементов этих троек.
  \filehint{17.txt}
  \answer{Два числа}
  \solution{Предварительно вычисляем $\min\_r5 = \min\_elem \% 5$ и $\max\_r7 = \max\_elem \% 7$. Перебираем тройки, проверяем все три условия.}

\end{enumerate}

\bigskip

% ========================
%  СЛОЖНЫЕ ЗАДАЧИ (hard)
% ========================

\subsection*{\textcolor{hardcolor}{Сложный уровень}}

\begin{enumerate}[label=\textbf{Задача \arabic*.}, leftmargin=*, itemsep=10mm]

  \item % ID: cs_ege_17_hard_001 | \hardtag | src: №75281
  \hardtag\ Файл содержит последовательность натуральных чисел, не превышающих $100\,000$.
  Назовём тройкой три \emph{идущих подряд} элемента последовательности.
  Определите количество троек, для которых выполняются следующие условия:
  \begin{itemize}
    \item в тройке есть четырёхзначные числа;
    \item в тройке \textbf{не более одного} числа, у которого остаток от деления
    на 5 равен остатку от деления на 5 \emph{максимального} элемента
    всей последовательности;
    \item в тройке \textbf{не менее двух} чисел, у которых остаток от деления на 7
    равен остатку от деления на 7 \emph{минимального} элемента всей
    последовательности.
  \end{itemize}
  В ответе запишите два числа: сначала количество найденных троек, затем
  максимальную величину суммы элементов этих троек.
  \filehint{17.txt}
  \answer{Два числа}
  \solution{Аналогично med\_005, но остатки берём от максимума (для /5) и минимума (для /7).}

  \bigskip

  \item % ID: cs_ege_17_hard_002 | \hardtag | src: №76120
  \hardtag\ В файле \texttt{17-1.txt} содержится последовательность целых чисел.
  Все элементы различны и могут принимать значения от $-100\,000$ до $100\,000$.
  Определите количество троек последовательности, у которых:
  \begin{itemize}
    \item \textbf{старшие разряды чисел совпадают};
    \item \textbf{хотя бы одно} из чисел оканчивается на 7 и является трёхзначным;
    \item \textbf{модуль суммы троек} меньше максимального элемента
    последовательности, оканчивающегося на 7.
  \end{itemize}
  В ответе запишите количество найденных троек, затем модуль максимальной из сумм
  элементов таких троек. Под тройкой подразумевается три \emph{идущих подряд}
  элемента.
  \filehint{17-1.txt}
  \answer{Два числа}
  \solution{«Старший разряд» — первая цифра числа (для отрицательных — первая цифра модуля). Предварительно находим нужный максимум. Перебираем тройки, проверяем все условия.}

  \bigskip

  \item % ID: cs_ege_17_hard_003 | \hardtag | src: №76685
  \hardtag\ Файл содержит последовательность натуральных чисел, не превышающих $100\,000$.
  Назовём тройкой три \emph{идущих подряд} элемента последовательности.
  Определите количество троек, для которых выполняются следующие условия:
  \begin{itemize}
    \item в тройке есть хотя бы \textbf{два четырёхзначных числа};
    \item в тройке есть число, \textbf{последняя цифра которого совпадает с
    последней цифрой максимального элемента} всей последовательности;
    \item в тройке \textbf{нет чисел}, последняя цифра которых совпадает с
    последней цифрой \emph{минимального} элемента всей последовательности.
  \end{itemize}
  В ответе запишите два числа: сначала количество найденных троек, затем
  максимальную величину суммы элементов этих троек.
  \filehint{17.txt}
  \answer{Два числа}
  \solution{Предварительно вычисляем $\max\_d = \max\_elem \% 10$ и $\min\_d = \min\_elem \% 10$. Условие: $\geq 2$ четырёхзначных, хотя бы одно $\% 10 = \max\_d$, ни одно $\%10 \neq \min\_d$.}

  \bigskip

  \item % ID: cs_ege_17_hard_004 | \hardtag | src: №64902
  \hardtag\ Файл содержит последовательность натуральных чисел, не превышающих $100\,000$.
  Назовём тройкой три \emph{идущих подряд} элемента последовательности.
  Определите количество троек, для которых выполняются следующие условия:
  \begin{itemize}
    \item в тройке есть \textbf{пятизначные числа, но не все числа в тройке
    пятизначные};
    \item в тройке \textbf{больше чисел, кратных 3}, чем чисел, кратных 5;
    \item \textbf{сумма элементов тройки больше максимального элемента}
    последовательности, запись которого заканчивается на $238$.
  \end{itemize}
  Гарантируется, что в последовательности есть хотя бы один такой элемент.
  В ответе запишите два числа: сначала количество найденных троек, затем
  максимальную величину суммы элементов этих троек.
  \filehint{17.txt}
  \answer{Два числа}
  \solution{Предварительно находим $\max238$ — максимум среди чисел, оканчивающихся на 238. Перебираем тройки, проверяем три условия.}

  \bigskip

  \item % ID: cs_ege_17_hard_005 | \hardtag | src: №61397
  \hardtag\ Файл содержит последовательность натуральных чисел, не превышающих $100\,000$.
  Назовём тройкой три \emph{идущих подряд} элемента последовательности.
  Определите количество троек, для которых выполняются следующие условия:
  \begin{itemize}
    \item \textbf{ровно два числа} в тройке четырёхзначные;
    \item \textbf{хотя бы одно число} в тройке делится на 5;
    \item \textbf{сумма элементов тройки больше максимального элемента}
    последовательности, запись которого заканчивается на $17$.
  \end{itemize}
  Гарантируется, что в последовательности есть хотя бы один такой элемент.
  В ответе запишите два числа: сначала количество найденных троек, затем
  максимальную величину суммы элементов этих троек.
  \filehint{17.txt}
  \answer{Два числа}
  \solution{Предварительно находим $\max17$ — максимум среди чисел, оканчивающихся на 17. Ровно два четырёхзначных: считаем количество четырёхзначных в тройке = 2. Проверяем делимость на 5 и порог суммы.}

\end{enumerate}

\end{document}
